\par\noindent\begin{minipage}{\linewidth}
% !TeX root = ../教材题库_全部条目_转写工作稿.tex
% Source: XJ-B1-C03-TH-P079-Y0868
% 第3章 / 3.2 函数的基本性质 / 教材 p.79 / PDF 第86页
% 已对照原页转写题干及公式；待教师复核；未制作答案
% 原题及相邻原解截图（核对用）：../题目原图/XJ-B1-C03-TH-P079-Y0868-01.webp
% 本文件只在转写完成后作为题干使用；OCR 条目尚不是可靠题干。
\begin{bankitem}{p.79 思考}{第3章 / 3.2 函数的基本性质 / 教材 p.79 / PDF 第86页}{XJ-B1-C03-TH-P079-Y0868}{已对照原页转写题干及公式；待教师复核；未制作答案}
图 3.2-1 是湖北省武汉市某年 11 月 22 日 16 时到 23 日 16 时这 24 个小时内的气温变化图。该曲线给出了时间与气温间的函数关系。你能从中看出它的变化趋势和最高（或最低）温度吗？
\bankfig{C03-TH079}{0.78}
\end{bankitem}
\end{minipage}\par\medskip
